\section*{Capítulo \ref{ch:g10:soil-formation} --- El suelo, un recurso de formación lenta}

\begin{solution}{exo:g10:soil-formation:1}
La \omterm{def:g10:soil-formation:chemical-weathering}{meteorización química} es la transformación de los minerales de una roca
por reacciones con el agua y con los gases que lleva (dióxido de carbono,
oxígeno): algunos minerales se disuelven, otros se convierten en minerales
nuevos, estables en la superficie. La \omterm{def:g10:soil-formation:hydrolysis}{hidrólisis} es la \omterm{def:g10:soil-formation:chemical-weathering}{meteorización química}
de un silicato por el agua y sus iones hidrógeno: los iones metálicos pasan a
la disolución y el resto del mineral se reorganiza en un \omterm{def:g10:soil-formation:clay}{mineral de la
arcilla}.
\end{solution}

\begin{solution}{exo:g10:soil-formation:2}
La \omterm{def:g10:soil-formation:weathering-profile}{saprolita} es roca descompuesta en su sitio: blanda y terrosa, conserva los
granos, las bandas y la estructura de su \omterm{def:g10:soil-formation:weathering-profile}{roca madre}, y pasa hacia abajo, sin
ruptura, a la roca sana. Un suelo formado sobre arena de río descansa sobre un
material traído de otro lugar: sus granos están redondeados y seleccionados,
y no muestra nada de la estructura de la roca que tiene debajo.
\end{solution}

\begin{solution}{exo:g10:soil-formation:3}
Los feldespatos se convierten en arcilla (caolinita) y liberan al agua
potasio, sodio o calcio, y sílice; la mica negra se convierte en arcilla y en
óxidos de hierro; el cuarzo queda intacto, en forma de arena.
\end{solution}

\begin{solution}{exo:g10:soil-formation:4}
El hierro de los minerales oscuros, \ce{Fe^2+}, es oxidado por el oxígeno que
lleva el agua y forma óxidos de hierro(III) insolubles: la goethita, de color
pardo amarillento, y la hematites, roja, que recubren los granos y tiñen todo
el material.
\end{solution}

\begin{solution}{exo:g10:soil-formation:5}
\qty{0.017}{mm} al año y \qty{1.537}{mm} al año: $1.537/0.017 \approx 90$
veces mayor.
\end{solution}

\begin{solution}{exo:g10:soil-formation:6}
Calcio: 1 y 1. Carbono: $1 + 1 = 2$ y 2. Oxígeno: $3 + 2 + 1 = 6$ y
$2 \times 3 = 6$. Hidrógeno: 2 y 2. Carga: 0, y $+2 - 2 = 0$.
\end{solution}

\begin{solution}{exo:g10:soil-formation:7}
La disolución de la calcita solo produce iones disueltos, \ce{Ca^2+} y
\ce{HCO3-}: todo se va con el agua, y una caliza pura no deja ningún residuo
sólido. La \omterm{def:g10:soil-formation:hydrolysis}{hidrólisis} del feldespato produce un sólido, la caolinita, que se
queda; y el cuarzo del granito apenas reacciona. Por eso el granito deja
arcilla y arena.
\end{solution}

\begin{solution}{exo:g10:soil-formation:8}
Las escamas de arcilla y el humus llevan cargas negativas que retienen en sus
superficies los iones potasio positivos \ce{K+}, de modo que el agua de
lluvia que pasa no se los lleva. Una raíz libera iones hidrógeno, que ocupan
el lugar del potasio en las escamas; la raíz absorbe el potasio liberado en
el agua del suelo.
\end{solution}

\begin{solution}{exo:g10:soil-formation:9}
La \omterm{def:g10:soil-formation:chemical-weathering}{meteorización química} necesita agua, abundante en los climas húmedos y
escasa en los secos (y helada en los fríos); y sus reacciones van más deprisa
cuando hace calor. Solo con calor y agua abundante, durante mucho tiempo,
llegan las reacciones lo bastante lejos como para eliminar casi todo lo
soluble y dejar óxidos de hierro y de aluminio.
\end{solution}

\begin{solution}{exo:g10:soil-formation:10}
$10/0.017 \approx 590$ años para \qty{1}{cm}; $300/0.017 \approx 17\,600$
años para \qty{30}{cm}; $1000/0.017 \approx 59\,000$ años para \qty{1}{m}.
\end{solution}

\begin{solution}{exo:g10:soil-formation:11}
Volumen: $10\,000 \times 0.0015 = \qty{15}{m^3}$; masa: $15 \times 1200 =
\qty{18000}{kg}$, \num{18}~toneladas por hectárea y por año.
\end{solution}

\begin{solution}{exo:g10:soil-formation:12}
$46\,000/0.017 \approx 2.7$ millones de años. Una \omterm{def:g10:soil-formation:weathering-profile}{saprolita} tan gruesa
necesita un paisaje que haya permanecido estable, sin que le arranquen su
manto meteorizado, durante millones de años.
\end{solution}

\begin{solution}{exo:g10:soil-formation:13}
Goethita: \ce{4 Fe^2+ + O2 + 6 H2O -> 4 FeO(OH) + 8 H+}: hierro 4 y 4,
oxígeno $2 + 6 = 8$ y $4 \times 2 = 8$, hidrógeno 12 y $4 + 8 = 12$, carga
$+8$ y $+8$. Hematites: \ce{4 Fe^2+ + O2 + 4 H2O -> 2 Fe2O3 + 8 H+}: hierro 4
y 4, oxígeno $2 + 4 = 6$ y 6, hidrógeno 8 y 8, carga $+8$ y $+8$.
\end{solution}

\begin{solution}{exo:g10:soil-formation:14}
Pérdida neta: $0.082 - 0.017 = \qty{0.065}{mm}$ al año; $300/0.065 \approx 4600$ años. Con el arado convencional, $300/1.52 \approx 200$ años: más de
veinte veces menos.
\end{solution}

\begin{solution}{exo:g10:soil-formation:15}
La tonelada contiene \qty{300}{kg} de ortosa. Cada kilogramo da
\qty{464}{g} de caolinita: $300 \times 0.464 \approx \qty{139}{kg}$. Potasio:
$300 \times 0.1405 \approx \qty{42}{kg}$.
\end{solution}

\begin{solution}{pb:g10:soil-formation:1}
\textbf{1.} $4 \times 0.3048 \approx \qty{1.22}{m}$.

\textbf{2.} $0.5 \times 30.48 \approx \qty{15}{cm}$.

\textbf{3.} El humus, la materia orgánica oscura que los descomponedores
fabrican a partir de las plantas muertas.

\textbf{4.} Los óxidos de hierro, goethita y hematites, formados por la
oxidación del hierro(II) de los minerales del material original.

\textbf{5.} No. Por debajo, el perfil pasaría a roca meteorizada, la
\omterm{def:g10:soil-formation:weathering-profile}{saprolita} (u otro material meteorizado), y después a la \omterm{def:g10:soil-formation:weathering-profile}{roca madre} sana.

\textbf{6.} \ce{2 KAlSi3O8 + 2 H+ + 9 H2O -> Al2Si2O5(OH)4 + 2 K+ + 4 H4SiO4}.

\textbf{7.} $1000/278.33 \approx \qty{3.59}{mol}$.

\textbf{8.} $3.59/2 \approx \qty{1.80}{mol}$ de caolinita, $1.80 \times 258.16
\approx \qty{464}{g}$.

\textbf{9.} Los iones potasio y la sílice disuelta: el agua los lleva a los
ríos, y los ríos, al mar.

\textbf{10.} Apenas reacciona: los granos de cuarzo quedan como arena.

\textbf{11.} $1219/0.017 \approx 72\,000$ años.

\textbf{12.} $1.537 - 0.017 = \qty{1.52}{mm}$ al año.

\textbf{13.} $1219/1.52 \approx 800$ años.

\textbf{14.} $72\,000/800 = 90$: el perfil se pierde noventa veces más
deprisa de lo que se formó.

\textbf{15.} No: $0.013 < 0.017$; la erosión es algo más lenta que la
producción, y el suelo se mantiene.

\textbf{16.} $T = S/(E - P)$.

\textbf{17.} $300/1.52 \approx 200$ años.

\textbf{18.} $300/(0.082 - 0.017) = 300/0.065 \approx 4600$ años.

\textbf{19.} Un suelo perdido en unos pocos siglos tarda decenas de miles de
años en volver a formarse: a la escala de las vidas humanas y de las
sociedades, lo que se pierde no se repone.

\textbf{20.} $1000/1.52 \approx 660$ años: \textbf{con el arado convencional,
un metro de suelo se pierde en unos 660 años, unos pocos siglos}.
\end{solution}
